% Integrated after independent root, criteria, and geometric audits, 2026-11-16.
% Original author draft retained in agent_F/F_fixed_r_configuration_draft.tex.
\subsection{A fixed-prime repair of configuration fiber exactness}

\begin{claim}{AG.41}{Proved: fixed-prime cuspidal configuration exactness}
\label{clm:AGfixedrconfiguration}
Let \(p\) and \(r\) be distinct primes. Let \(C\) be a
smooth proper connected complex curve of genus at least two,
let \(x\in C\), and set
\[
 U=(C\times C)-\Delta_C,\qquad V=C-\{x\},\qquad
 \Delta=\pi_1(C)^{(p')},\quad F=\pi_1(V)^{(p')},\quad
 E=\pi_1(U)^{(p')}.
                                                        \tag{AG.41.1}
\]
The fiber embedding uses the second projection \(U\to C\).
Define
\[
\begin{aligned}
 K_F&=\ker(F\to\Delta),
 &K_E&=\ker(E\to\Delta\times\Delta),\\
 F_{\mathrm{c}r}&=F/\ker(K_F\to K_F^{(r)}),
 &E_{\mathrm{c}r}&=E/\ker(K_E\to K_E^{(r)}).
\end{aligned}
                                                        \tag{AG.41.2}
\]
Then the fiber map induces an exact sequence
\[
 1\longrightarrow F_{\mathrm{c}r}
 \longrightarrow E_{\mathrm{c}r}
 \xrightarrow{\mathrm{pr}_2}\Delta\longrightarrow1.
                                                        \tag{AG.41.3}
\]
The same statement holds for geometric etale prime-to-\(p\)
groups over every algebraically closed field of characteristic
\(p\), with the corresponding characteristic relative quotients.

More precisely, every finite quotient of \(F_{\mathrm{c}r}\)
is dominated, on the fiber, by a finite quotient of
\(E_{\mathrm{c}r}\). The finite quotient of the total group
can be chosen to have a normal \(r\)-subgroup whose quotient
is a finite prime-to-\(p\) quotient of
\(\Delta\times\Delta\).
\end{claim}

The distinction from full prime-to-\(p\) exactness is essential.
Here the entire cuspidal kernel, on both sides, is replaced by
its maximal pro-\(r\) quotient for one fixed prime \(r\).
No injection of \(F\) into \(E\) is assumed or concluded.

\begin{proof}
\emph{1. A pro-\(r\) monodromy lemma without a section.}
Let \(1\to L\to H\to B\to1\) be an exact sequence of
discrete groups, where \(L\) is a free group of finite rank
at least two. Suppose the image of conjugation on
\(L/[L,L]L^r\) is an \(r\)-group. Then
\[
                         L^{(r)}\longrightarrow H^{(r)}
                    \quad\hbox{is injective}.             \tag{AG.41.4}
\]
Indeed, conjugation by \(H\) acts canonically on the
characteristic completion \(L^{(r)}\). The kernel of
\[
 \operatorname{Aut}(L^{(r)})\longrightarrow
 \operatorname{Aut}(L^{(r)}/\Phi(L^{(r)}))
                                                        \tag{AG.41.5}
\]
is pro-\(r\). Here is the finite argument, including the
topological passage. For a finite \(r\)-group \(P\), fix
an ordered basis of \(P/\Phi(P)\), of size \(d\).
There are exactly \(|\Phi(P)|^d\) tuples lifting that basis,
and every such tuple generates \(P\), by the elementary
Frattini criterion. The kernel of the automorphism action
on \(P/\Phi(P)\) acts freely on these tuples: an
automorphism fixing a generating tuple is the identity.
Its order therefore divides \(|\Phi(P)|^d\), so it is an
\(r\)-group. The iterated Frattini subgroups of the
finitely generated pro-\(r\) group \(L^{(r)}\) are
characteristic, open, and cofinal: their images in every
finite \(r\)-quotient eventually vanish. Applying the
finite argument to these quotients proves (AG.41.5).

It follows that the closure of the conjugation image of
\(H\) in \(\operatorname{Aut}(L^{(r)})\) is pro-\(r\).
The homomorphism from \(H\) to that closure extends to
\(H^{(r)}\). On \(L^{(r)}\) its composite is the inner
automorphism map. A nonabelian free pro-\(r\) group has
trivial center, so this composite is injective; see, for
example, \cite[Theorem~5.1, proof Case~I, and
Corollary~5.3]{FF_KochloukovaZalesskii2011}.
This proves (AG.41.4). In particular no splitting of
\(H\to B\), topological section, or lift of an outer
action is needed.

The compact embedding in (AG.41.4) induces the full
pro-\(r\) topology on the fiber. Explicitly, if
\(W\triangleleft L^{(r)}\) is open, there is an open
normal subgroup \(N\triangleleft H^{(r)}\) with
\(N\cap L^{(r)}\subset W\). Thus the finite quotient
\(H^{(r)}/N\) dominates the prescribed quotient of the
fiber; equality with that quotient is unnecessary.

\emph{2. Finite proper covers make the monodromy an
\(r\)-group modulo \(r\).}
Let \(a:C'\to C\) be a connected finite Galois etale
cover with group \(A\) of order prime to \(p\). Form
\[
 Y=(C'\times C')-(a\times a)^{-1}(\Delta_C).
                                                        \tag{AG.41.6}
\]
The map \(Y\to U\) is Galois with group \(A\times A\).
The second projection \(Y\to C'\) has fiber
\(V'=C'-a^{-1}(x)\), after choosing a point over \(x\)
in the second factor. Its removed divisor is the disjoint
union of the graphs of all deck transformations of \(a\).
Thus every puncture is labeled throughout this family.
The topological fibration is locally trivial; since the
closed surface \(C'\) is aspherical, its homotopy sequence
gives
\[
 1\longrightarrow\pi_1(V')\longrightarrow\pi_1(Y)
 \longrightarrow\pi_1(C')\longrightarrow1.
                                                        \tag{AG.41.7}
\]
The fiber group is free of rank
\(2g(C')+|A|-1\), in particular of rank at least two.

In homology with coefficients in \(\mathbb F_r\), the
boundary loops span a subspace \(W\) and give an exact
sequence
\[
 0\longrightarrow W\longrightarrow H_1(V',\mathbb F_r)
 \longrightarrow H_1(C',\mathbb F_r)\longrightarrow0.
                                                        \tag{AG.41.8}
\]
The boundary labels and complex orientations make monodromy
the identity on \(W\). The compactified family is the
product \(C'\times C'\to C'\), so monodromy is also
the identity on the last term. Therefore its image lies
in the additive shear group
\(\operatorname{Hom}_{\mathbb F_r}(H_1(C',\mathbb F_r),W)\),
a finite \(r\)-group. Fiber conjugations already act
trivially on homology. Step~1 consequently proves
\[
             \pi_1(V')^{(r)}\hookrightarrow\pi_1(Y)^{(r)}.
                                                        \tag{AG.41.9}
\]

\emph{3. Every finite relative fiber quotient is detected.}
Take a continuous surjection
\(q:F_{\mathrm{c}r}\to Q\), with \(Q\) finite, and
put \(R=q(K_F^{(r)})\). This is a normal \(r\)-subgroup
of \(Q\); the quotient \(A=Q/R\) is a finite quotient
of \(\Delta\). Use its Galois proper cover \(C'\to C\)
in Step~2. On discrete fundamental groups write
\(\Gamma=\pi_1(U)\), \(\Gamma_0=\pi_1(Y)\), and
\(\Lambda=\pi_1(V)\), \(\Lambda_0=\pi_1(V')\).
Then \(\Gamma_0\triangleleft\Gamma\),
\(\Gamma/\Gamma_0=A\times A\), and
\(\Lambda\cap\Gamma_0=\Lambda_0\).
The restriction of \(q\) to \(\Lambda_0\) has image
in \(R\), so factors through a finite \(r\)-quotient
of \(\Lambda_0^{(r)}\).

By (AG.41.9) and the final paragraph of Step~1 there is
a finite \(r\)-quotient of \(\Gamma_0\), with kernel
\(J\triangleleft\Gamma_0\), such that
\[
                         J\cap\Lambda_0\subset\ker(q).
                                                        \tag{AG.41.10}
\]
Take its normal core in \(\Gamma\):
\[
                  J^*=\bigcap_{g\in\Gamma/\Gamma_0}gJg^{-1}.
                                                        \tag{AG.41.11}
\]
This finite intersection is legitimate because \(J\) is
normal in \(\Gamma_0\). The quotient
\(\Gamma_0/J^*\) embeds in a finite product of
\(r\)-groups, hence is an \(r\)-group. Consequently
\(\Gamma/J^*\) is a finite prime-to-\(p\) group with
normal \(r\)-subgroup \(\Gamma_0/J^*\) and quotient
\(A\times A\). Moreover
\[
                 J^*\cap\Lambda\subset\ker(q),
                                                        \tag{AG.41.12}
\]
because \(J^*\subset\Gamma_0\) and (AG.41.10) applies.
Thus this total-space quotient dominates \(q\) on the
fiber. The inclusion of finite kernels continues to hold
after prime-to-\(p\) completion: it can be checked in
the joint finite image of the two quotient maps, in which
the discrete group has dense, hence surjective, image.

The quotient \(\Gamma\to\Gamma/J^*\) factors through
\(E_{\mathrm{c}r}\), not merely through \(E\).
Indeed its image of \(K_E\) lies in
\(\Gamma_0/J^*\), since \(\Gamma/\Gamma_0\) comes
from the proper product cover. That image is therefore an
\(r\)-group, and the kernel of \(K_E\to K_E^{(r)}\)
is killed. This proves the claimed cofinal finite-quotient
statement.

\emph{4. Injectivity and exactness.}
The fiber map factors through \(F_{\mathrm{c}r}\),
because it takes \(K_F\) to the pro-\(r\) kernel in
\(E_{\mathrm{c}r}\). Every finite quotient of that fiber
group is dominated by one from Step~3, so its kernel is
trivial. This proves the first arrow of (AG.41.3).

The ordinary topological fibration
\(C-\{x\}\to U\to C\) is right exact on fundamental
groups. Right exactness survives maximal prime-to-\(p\)
completion: quotienting \(E\) by the closed image of
the fiber has exactly the finite prime-to-\(p\) quotients
of \(\pi_1(C)\), hence is \(\Delta\). The subgroup
killed to form \(E_{\mathrm{c}r}\) is contained in
this fiber kernel. The image of the compact group
\(F_{\mathrm{c}r}\) is closed and is therefore exactly
the kernel of \(\mathrm{pr}_2\). This proves (AG.41.3).

\emph{5. Characteristic-\(p\) specialization.}
Spread out the pointed curve over a finitely generated
subfield, pass to a countable algebraic closure, and lift
the smooth pointed proper curve to a mixed-characteristic
complete discrete valuation ring. Curve deformations are
unobstructed, and proper formal lifts algebraize. An
algebraic closure of the fraction field embeds abstractly
in \(\mathbb C\). Both
\((\mathcal C,\mathcal x)\) and
\((\mathcal C\times\mathcal C,\Delta_{\mathcal C})\)
are smooth proper pairs with relative normal-crossing
boundary. SGA~1 XIII, Corollary~2.9 and Section~2.10 give
functorial isomorphisms of their geometric prime-to-\(p\)
groups \cite{AR_SGA1}; this is the same comparison used
in Claim~\ref{clm:AGconfigurationodd}.
These isomorphisms identify \(F,E,\Delta\), their proper
quotients and the fiber maps. They consequently identify
the characteristic relative quotients (AG.41.2).
The complex result therefore proves the asserted
characteristic-\(p\) result.
\end{proof}

\paragraph{Arithmetic consequence, with the section specified.}
Let \(C/k\) instead be a smooth proper geometrically
connected curve of genus at least two over a finite field
of characteristic \(p\). Write \(\Pi_C\) for its
arithmetic geometrically prime-to-\(p\) group and let
\(\mathcal E\) be the corresponding group of
\(U=(C\times C)-\Delta_C\). Form its cuspidally
pro-\(r\) quotient relative to
\[
                \Pi_C\times_{G_k}\Pi_C.
                                                        \tag{AG.41.13}
\]
The kernel relative to this product is geometric, so this
quotient has geometric subgroup \(E_{\mathrm{c}r}\)
from (AG.41.2). If \(P\in C(k)\) and
\(s_P:G_k\to\Pi_C\) is its actual point section,
the natural map from the arithmetic cuspidally pro-\(r\)
group of \(C-P\) identifies it with
\[
       \mathcal E_{\mathrm{c}r}
                  \times_{\Pi_C,\mathrm{pr}_2,s_P}G_k.
                                                        \tag{AG.41.14}
\]
Indeed on geometric kernels this is (AG.41.3), and on
arithmetic quotients it is the identity of \(G_k\).
The elementary extension comparison therefore makes it
an isomorphism. Since a section is injective, the right
side can equivalently be viewed as the inverse image of
\(s_P(G_k)\) under the second projection. Choices of
basepoint paths give the usual geometric inner ambiguity;
no section of the configuration fibration was postulated.

\paragraph{The boundary comparison.}
Claim~\ref{clm:AGcuspidaltransport} identifies the actual
puncture decomposition group with the intersection with
the global diagonal boundary decomposition group. This
is a separate tame local calculation; it is not silently
included in the geometric fiber theorem.
